\documentclass{jlreq}
\usepackage{amsmath,amssymb,hyperref}
\usepackage[thmtools-compat]{keytheorems}
\declaretheoremstyle[headfont=\normalfont\bfseries, bodyfont=\normalfont, headpunct={}, postheadspace={ }, spaceabove=3pt, spacebelow=3pt, qed=$\square$]{boxed}
\declaretheorem[name=定義, style=boxed, within=section]{definition}
\declaretheorem[name=定理, style=boxed, sibling=definition]{theorem}
\declaretheorem[name=注意, style=boxed, sibling=definition]{remark}
\renewcommand{\proofname}{\bfseries 証明}
\renewcommand{\qedsymbol}{$\blacksquare$}

\begin{document}
定義~\ref{split}, 定義~\ref{after} を参照.
\section{圏}
本文の段落です. これは長い文章でページを埋めるためのものです. 本文の段落です. これは長い文章でページを埋めるためのものです. 本文の段落です. これは長い文章でページを埋めるためのものです. 本文の段落です. これは長い文章でページを埋めるためのものです. 本文の段落です. これは長い文章でページを埋めるためのものです. 本文の段落です. これは長い文章でページを埋めるためのものです. 本文の段落です. これは長い文章でページを埋めるためのものです. 本文の段落です. これは長い文章でページを埋めるためのものです. 本文の段落です. これは長い文章でページを埋めるためのものです. 本文の段落です. これは長い文章でページを埋めるためのものです. 本文の段落です. これは長い文章でページを埋めるためのものです. 本文の段落です. これは長い文章でページを埋めるためのものです. 本文の段落です. これは長い文章でページを埋めるためのものです. 本文の段落です. これは長い文章でページを埋めるためのものです. 本文の段落です. これは長い文章でページを埋めるためのものです. 本文の段落です. これは長い文章でページを埋めるためのものです. 本文の段落です. これは長い文章でページを埋めるためのものです. 本文の段落です. これは長い文章でページを埋めるためのものです. 本文の段落です. これは長い文章でページを埋めるためのものです. 本文の段落です. これは長い文章でページを埋めるためのものです. 本文の段落です. これは長い文章でページを埋めるためのものです. 本文の段落です. これは長い文章でページを埋めるためのものです. 本文の段落です. これは長い文章でページを埋めるためのものです. 本文の段落です. これは長い文章でページを埋めるためのものです. 本文の段落です. これは長い文章でページを埋めるためのものです. 本文の段落です. これは長い文章でページを埋めるためのものです. 本文の段落です. これは長い文章でページを埋めるためのものです. 本文の段落です. これは長い文章でページを埋めるためのものです. 本文の段落です. これは長い文章でページを埋めるためのものです. 本文の段落です. これは長い文章でページを埋めるためのものです. 本文の段落です. これは長い文章でページを埋めるためのものです. 本文の段落です. これは長い文章でページを埋めるためのものです. 本文の段落です. これは長い文章でページを埋めるためのものです. 本文の段落です. これは長い文章でページを埋めるためのものです. 本文の段落です. これは長い文章でページを埋めるためのものです. 本文の段落です. これは長い文章でページを埋めるためのものです. 本文の段落です. これは長い文章でページを埋めるためのものです. 本文の段落です. これは長い文章でページを埋めるためのものです. 本文の段落です. これは長い文章でページを埋めるためのものです. 本文の段落です. これは長い文章でページを埋めるためのものです. 本文の段落です. これは長い文章でページを埋めるためのものです. 本文の段落です. これは長い文章でページを埋めるためのものです. 本文の段落です. これは長い文章でページを埋めるためのものです. 本文の段落です. これは長い文章でページを埋めるためのものです. 本文の段落です. これは長い文章でページを埋めるためのものです. 本文の段落です. これは長い文章でページを埋めるためのものです. 本文の段落です. これは長い文章でページを埋めるためのものです. 本文の段落です. これは長い文章でページを埋めるためのものです. 本文の段落です. これは長い文章でページを埋めるためのものです. 本文の段落です. これは長い文章でページを埋めるためのものです. 本文の段落です. これは長い文章でページを埋めるためのものです. 本文の段落です. これは長い文章でページを埋めるためのものです. 本文の段落です. これは長い文章でページを埋めるためのものです. 本文の段落です. これは長い文章でページを埋めるためのものです. 本文の段落です. これは長い文章でページを埋めるためのものです. 本文の段落です. これは長い文章でページを埋めるためのものです. 本文の段落です. これは長い文章でページを埋めるためのものです. 本文の段落です. これは長い文章でページを埋めるためのものです. 本文の段落です. これは長い文章でページを埋めるためのものです. 本文の段落です. これは長い文章でページを埋めるためのものです. 本文の段落です. これは長い文章でページを埋めるためのものです. 本文の段落です. これは長い文章でページを埋めるためのものです. 本文の段落です. これは長い文章でページを埋めるためのものです. 本文の段落です. これは長い文章でページを埋めるためのものです. 本文の段落です. これは長い文章でページを埋めるためのものです. 本文の段落です. これは長い文章でページを埋めるためのものです. 本文の段落です. これは長い文章でページを埋めるためのものです. 本文の段落です. これは長い文章でページを埋めるためのものです. 本文の段落です. これは長い文章でページを埋めるためのものです. 本文の段落です. これは長い文章でページを埋めるためのものです. 本文の段落です. これは長い文章でページを埋めるためのものです. 本文の段落です. これは長い文章でページを埋めるためのものです. 本文の段落です. これは長い文章でページを埋めるためのものです. 本文の段落です. これは長い文章でページを埋めるためのものです. 本文の段落です. これは長い文章でページを埋めるためのものです. 本文の段落です. これは長い文章でページを埋めるためのものです. 本文の段落です. これは長い文章でページを埋めるためのものです. 本文の段落です. これは長い文章でページを埋めるためのものです. 本文の段落です. これは長い文章でページを埋めるためのものです. 本文の段落です. これは長い文章でページを埋めるためのものです. 

\begin{definition}\label{split}
$C$を局所小圏, $F \colon C \to \mathbf{Set}$を共変関手とする. 次が成り立つとき$u$は$F$の普遍元であるという.
\[ \forall B \in \mathrm{Ob}(C), \forall x \in F(B), \exists! f \in \mathrm{Mor}(A, B), F(f)(u) = x \tag{$\ast$} \]
このとき組$(A,u)$は$F$の普遍対であるという. これは定義の続きの文章であり長く書かれている. このとき組$(A,u)$は$F$の普遍対であるという. これは定義の続きの文章であり長く書かれている. このとき組$(A,u)$は$F$の普遍対であるという. これは定義の続きの文章であり長く書かれている. このとき組$(A,u)$は$F$の普遍対であるという. これは定義の続きの文章であり長く書かれている. このとき組$(A,u)$は$F$の普遍対であるという. これは定義の続きの文章であり長く書かれている. このとき組$(A,u)$は$F$の普遍対であるという. これは定義の続きの文章であり長く書かれている. 
\begin{itemize}
\item 対象の集まり$\mathrm{Ob}(C)$.
\item 各対象の組$A, B$に対して射の集合$\mathrm{Mor}(A,B)$.
\end{itemize}
これらが結合律と単位律を満たすとき$C$を圏という.
\end{definition}
\begin{definition}\label{after}
次の定義である.
\end{definition}
本文の段落がここから始まる. これは定義の外である.
\end{document}
